\section{Dokumentation der KI-Nutzung}

Generative KI-Werkzeuge wurden bei der Analyse der Repositorystruktur, bei der
Planung und Überarbeitung von Quellcode, bei der Auswertung technischer Fehler
sowie bei der Strukturierung und sprachlichen Überarbeitung der Projektarbeit
unterstützend eingesetzt. Die fachliche Auswahl, Prüfung und Interpretation der
Ergebnisse verbleibt beim Verfasser.

\begin{table}[H]
\centering
\small
\begin{tabularx}{\textwidth}{p{3.0cm}p{2.8cm}p{2.5cm}X}
\toprule
Werkzeug/Modell & Anbieter & Zeitraum & Einsatzbereich\\
\midrule
GitHub Copilot, GPT-5.6 Luna & GitHub & August 2026 &
Repositoryanalyse, Implementierungsunterstützung, Fehlerdiagnose,
Versuchsplanung, Strukturierung und LaTeX-Überarbeitung.\\
\bottomrule
\end{tabularx}
\caption{Aktuell dokumentierter KI-Einsatz. Vor Abgabe sind alle tatsächlich
verwendeten Werkzeuge und Zeiträume zu ergänzen.}
\end{table}

Relevante Prompts, Antworten und dateibasierte Ausgaben sind in einem separaten
digitalen Nachweis zu gruppieren. Empfohlen werden mindestens die Gruppen
\emph{Pipelinekonzeption und Implementierung}, \emph{Fehleranalyse und Tests},
\emph{wissenschaftliche Auswertung} sowie \emph{LaTeX- und Sprachüberarbeitung}.
KI-Ausgaben werden nicht als wissenschaftliche Quellen zitiert. Fachliche
Aussagen sind anhand von Primärliteratur, ausführbarem Code und archivierten
Versuchsergebnissen zu prüfen.

\textbf{Offen vor Abgabe:} Zulässigkeit und gewünschte Granularität der
KI-Dokumentation sind mit dem Betreuer beziehungsweise Prüfer abzustimmen. Die
vollständige Eigenständigkeitserklärung der THWS ist unverändert zu übernehmen.
